\section{An explicit Bondi-radiative example and a detector limitation}
\label{sec:radiativeexample}
We now connect one coefficient functional to a conventional gravitational
soft contribution. This is a construction of smooth radiative data at
future null infinity, with an explicit memory and finite gravitational
flux. It is not a construction of a complete global scattering spacetime.
We use the unit round sphere metric $q_{AB}$ and the convention
$N_{AB}=\partial_u C_{AB}$ of \cite{pasterski}.

For a smooth real supertranslation test $F$, the soft part obtained
from equation~12 of \cite{pasterski} is
\begin{equation}\label{eq:bondisoft}
 Q_S[F]= -\frac{1}{16\pi G}\int_{\mathbb R\times S^2}
 F D_A D_B N^{AB}\,du\,d\Omega.
\end{equation}
Let $X$ be the Hilbert space of square-integrable symmetric trace-free
two-tensors on the sphere, with contractions taken using $q_{AB}$.
Writing
\[
 E_{AB}[F]=D_A D_BF-\tfrac12q_{AB}\Delta F,
 \qquad b_{AB}=\int_{\mathbb R}N_{AB}\,du,
\]
two integrations by parts on the closed sphere give
\begin{equation}\label{eq:bondifunctional}
 Q_S[F]=\ell_F(b),\qquad
 \ell_F(b)=-\frac{\langle E[F],b\rangle_X}{16\pi G},
 \qquad \|\ell_F\|=\frac{\|E[F]\|_X}{16\pi G}.
\end{equation}
For real $F$ this expression extends to a complex-linear functional on
the complexified tensor space. The norm equality follows by taking
$b$ aligned with $E[F]$ when this tensor is nonzero.

\begin{proposition}[A normalized harmonic charge]\label{prop:harmoniccharge}
Let $Y$ be a real scalar spherical harmonic with
$-\Delta Y=\lambda_lY$, $\lambda_l=l(l+1)$, and $\|Y\|_2=1$.
Then
\begin{equation}\label{eq:stfnorm}
 H_l:=\|E[Y]\|_X^2
 =\frac{\lambda_l(\lambda_l-2)}2
 =\frac{(l-1)l(l+1)(l+2)}2.
\end{equation}
For $l\ge2$, set $T_{AB}=E_{AB}[Y]/\sqrt{H_l}$ and $b=B T$,
where $B$ is real. Then
\begin{equation}\label{eq:harmoniccharge}
 \|b\|_X=|B|,\qquad Q_S[Y]= -\frac{B\sqrt{H_l}}{16\pi G}.
\end{equation}
A tensor-recovery error $\varepsilon_b$ in $X$ therefore gives a charge
error at most $\sqrt{H_l}\varepsilon_b/(16\pi G)$, with equality
possible for an aligned error. For $l=0,1$, the soft functional vanishes.
\end{proposition}
\begin{proof}
On the unit sphere $\operatorname{Ric}=q$. The integrated Bochner
identity yields
$\|D^2Y\|_2^2=\|\Delta Y\|_2^2-\|DY\|_2^2
=\lambda_l^2-\lambda_l$.
Removing the trace subtracts $\tfrac12\|\Delta Y\|_2^2$, proving
\eqref{eq:stfnorm}. Equations~\eqref{eq:bondifunctional} and
\eqref{eq:harmoniccharge} then give the assertions.
\end{proof}

There is no bound uniform over all $L^2$-normalized supertranslation
tests: $\|\ell_{Y_{lm}}\|=\sqrt{H_l}/(16\pi G)
\sim l^2/(16\pi G\sqrt2)$.
Thus a small $L^2$ error in memory does not give a common error constant
for arbitrarily fine angular tests. One may instead bound the tests
in $H^2(S^2)$. With the spectral convention
$\|F\|_{H^2}^2=\sum_{lm}(1+\lambda_l)^2|F_{lm}|^2$,
\[
 \|E[F]\|_X^2=\tfrac12\sum_{lm}\lambda_l(\lambda_l-2)|F_{lm}|^2
 \le\tfrac12\|F\|_{H^2}^2.
\]
This gives a uniform charge-transfer bound on an $H^2$-bounded set
of tests. Stronger angular regularity of the recovered memory is an
alternative route to controlling more singular dual observables.

\begin{proposition}[A finite-flux memory family]\label{prop:lorentziannews}
Fix a smooth real electric-parity tensor $b$ and $\tau>0$. Define
\begin{align}
 g_\tau(u)&=\frac{\tau}{\pi(u^2+\tau^2)},\qquad
 N_{AB}(u)=b_{AB}g_\tau(u),\label{eq:lorentziannews}\\
 C_{AB}(u)&=C^-_{AB}+b_{AB}
 \left[\frac12+\frac1\pi\arctan\frac u\tau\right].
\end{align}
Then $C^+-C^-=b$, and for the Fourier convention
$\widehat N(\omega)=\int e^{i\omega u}N(u)\,du$,
\begin{equation}\label{eq:newsfourier}
 \widehat N(\omega)=b e^{-\tau|\omega|},\qquad
 \int_{\mathbb R}\|N(u)\|_X^2du=\frac{\|b\|_X^2}{2\pi\tau}.
\end{equation}
In the convention with gravitational flux
$(32\pi G)^{-1}\int\|N\|_X^2du$, its radiated gravitational energy is
\begin{equation}\label{eq:radiatedenergy}
 E_{\rm grav}=\frac{\|b\|_X^2}{64\pi^2G\tau}.
\end{equation}
The positive-frequency news $f_\tau(\omega)=b e^{-\tau\omega}$ obeys
the $m=0$, $\alpha=1$ prior with $M=\tau\|b\|_X$ and
\begin{equation}\label{eq:newsmellin}
 G_\tau(z)=\Gamma(z)\tau^{-z}b,\qquad
 Q_a(f_\tau)=\|G_\tau\|_{a,\infty}
 =\Gamma(a)\tau^{-a}\|b\|_X,\quad a>0.
\end{equation}
Its Mellin residue at zero is $b$.
\end{proposition}
\begin{proof}
The elementary Lorentzian integrals give $\int g_\tau=1$ and
$\int g_\tau^2=1/(2\pi\tau)$. A contour integral at the pole
$u=i\tau$ for $\omega>0$, and at $u=-i\tau$ for $\omega<0$,
gives its Fourier transform; continuity gives the value at zero.
The gamma integral proves \eqref{eq:newsmellin}. The integral
triangle inequality bounds the whole vertical line by its zero-frequency
value, where equality is attained. Finally,
$\|f_\tau(\omega)-b\|_X
\le\tau\|b\|_X\omega$ and
$\|f'_\tau(0+)\|_X=\tau\|b\|_X$.
\end{proof}

For fixed nonzero $b$, letting $\tau\to\infty$ holds the memory and
its nonzero harmonic soft charge fixed while both the Mellin-line norm
and gravitational energy tend to zero. Thus a uniform upper bound on
radiated gravitational energy alone does not restore stability on this
radiative-data family. Its first frequency-derivative bound grows with
$\tau$, so the uniform prior in the recovery theorem excludes this
instability sequence. The theorem gives an actual bound for the
specified Bondi soft contribution by multiplication with the norm in
\eqref{eq:bondifunctional}.

For an unknown duration and a known uniform derivative bound $M>0$,
even this one-mode family has the same minimax exponent as the general
$\alpha=1$ theorem. To see this, fix a normalized tensor $T$ as in
Proposition~\ref{prop:harmoniccharge}. For $\eta>0$, take
\[
 \tau=\left(\frac{M\Gamma(a)}{\eta}\right)^{1/(a+1)},
 \qquad b_{\pm}=\pm\frac{M}{\tau}T.
\]
Both profiles have derivative norm $M$ and Mellin-line norm $\eta$.
The zero observation is compatible with both, so every estimator
has worst-case memory error at least $M/\tau$ and charge error at least
$\sqrt{H_l}M/(16\pi G\tau)$ for the corresponding harmonic test.
The upper bound from Theorem~\ref{thm:vectorlogs} with $m=0$ gives the
same order $M^{a/(a+1)}\eta^{1/(a+1)}$. In contrast, when the duration
$\tau$ is known exactly and the pulse shape is assumed exact, the
zero-frequency Mellin sample alone, $\tau^aY(0)/\Gamma(a)$, yields
memory error at most $\tau^a\eta/\Gamma(a)$. The opposite pair
$b_\pm=\pm\tau^a\eta T/\Gamma(a)$ shows this linear error bound is
optimal for the known-duration family with unrestricted amplitude.

\subsection{An idealized finite-time measurement}
Suppose an idealized measurement directly samples the news on
$[-U,U]$, with the pulse centred at zero. Its integrated memory is
\begin{align}
 b_U&=\frac2\pi\arctan(U/\tau)b,\label{eq:windowmemory}\\
 \|b-b_U\|_X&=\left[1-\frac2\pi\arctan(U/\tau)\right]\|b\|_X,\\
 \int_{-U}^U\|N(u)\|_X^2du
 &=\frac{\|b\|_X^2}{\pi^2}
 \left[\frac{U}{U^2+\tau^2}+\frac1\tau\arctan(U/\tau)\right].
 \label{eq:windowflux}
\end{align}
For fixed $U$ the measured fraction tends to zero as $\tau\to\infty$,
although the full memory remains $b$. If the additive measurement
error has $L^2([-U,U];X)$ norm at most $\epsilon$, integrating the noisy
record gives the explicit bound
\begin{equation}\label{eq:windowchargeerror}
 |\widehat Q_S[F]-Q_S[F]|
 \le \frac{\|E[F]\|_X}{16\pi G}
 \left\{\sqrt{2U}\epsilon+
 \left[1-\frac2\pi\arctan(U/\tau)\right]\|b\|_X\right\}.
\end{equation}
The first term is Cauchy--Schwarz and is attained by a temporally
constant error aligned with the angular test tensor. The second is
the exact missing memory for this pulse. The formulas follow by direct
integration; they do not derive a real instrument's strain response,
angular response, noise law, or quantum dressing prescription.

\subsection{Constraints, quantum states, and normalization}
Electric parity, for example $b=B E[Y]/\sqrt{H_l}$, avoids assuming
an arbitrary magnetic-parity vacuum transition. Even with this choice,
the other Bondi data must obey the constraints. In the conventions of
\cite{pasterski},
\[
 m_B^+-m_B^- =\tfrac14 D_A D_B b^{AB}-\int_{\mathbb R}T_{uu}\,du,
\]
where $T_{uu}$ includes the matter and quadratic gravitational terms.
The endpoint mass aspects, hard flux, and memory therefore cannot
all be prescribed independently. The construction above specifies
admissible smooth finite-flux radiative data and its conventional soft
functional; it proves neither a global Einstein evolution realizing
all chosen endpoints nor antipodal matching or the full scattering
Ward identity.

There is also an infrared representation issue. For nonzero $b$,
the ordinary zero-memory one-particle norm has an infrared contribution
proportional to
\[
 \int_0^\infty\frac{\|\widehat N(\omega)\|_X^2}{\omega}\,d\omega
 =\|b\|_X^2\int_0^\infty\frac{e^{-2\tau\omega}}{\omega}\,d\omega
 =\infty.
\]
Finite radiated energy does not put a nonzero-memory coherent
configuration into the ordinary zero-memory Fock representation;
the issue and memory representations are discussed in \cite{irfock}.
No dressed unitary scattering operator is supplied here.

Finally, $f_\tau$ is Fourier-transformed news, not an unstripped
graviton scattering amplitude. A conventional radiative mode expansion
relates news to an extra factor of $\omega$ times a shear mode, with
fixed polarization and gravitational normalization factors. If one
defines the stripped model $A(\omega)=f_\tau(\omega)/\omega$,
then $\Mell A(1+z)=G_\tau(z)$: the displaced line above is
$\Re\Delta=1+a$, not the boundary contour through the leading
soft pole. These distinctions delimit the explicit bridge established
here. They leave the complete physical scattering, dressing, and
all-orders questions open.
